\scanpage{91}
\diagram{\begin{tikzcd}[column sep=small,row sep=small]
&\Phi \arrow[r,dash]&N \arrow[ddr,dash]&\\
k \arrow[ur,"u"]\arrow[r,dash,"u'"]\arrow[dddr,"u''"']&\Phi' \arrow[ur,dash]\arrow[dddr,dash]&&\\
&\deleted{\unsure{\Phi'}}\arrow[ul,dash,"\deleted{\unsure{u'}}"]&&N''\\
&&&\\
&\Phi''\arrow[r,dash]&N'\arrow[uur,dash]&
\end{tikzcd}}

d'où
\diagram{\begin{tikzcd}
&\Phi\arrow[dr]&\\
k\arrow[ur]\arrow[r]\arrow[dr]&\Phi'\arrow[r]&N''\\
&\Phi''\arrow[ur]&
\end{tikzcd}}
\diagram{\begin{tikzcd}
\mathrm{DR}(\Phi)\arrow[dr,"\sim"]&\\
\mathrm{DR}(\Phi')\arrow[r,"\sim"]&\mathrm{DR}(N'')\\
\mathrm{DR}(\Phi'')\arrow[ur,"\sim"]&
\end{tikzcd}}

Argument analogue pour $f_*(\mathcal{DR}_{pd}^{**}(k,u))$, si $f:X\mathbin{-}Y$ morphisme de topos.

\uline{Condition} de \uline{De Rham sur Anneau} $k$ : la condition d'existence d'un $u:k\mathbin{-}\Phi$ tel que\ldots

Si $k$ satisfait condition de De Rham, de $=$ tte $k$-Algèbre. \unsure{Donc|Dans} le plus beau, c'est que $\mathbb Z_X$ y satisfasse ! C'est vrai pour

a) le topos ss $\widehat\Delta$

b) le topos défini par un espace topologique. \unsure{Voir cours ant.}

\uline{Question} Est-ce vrai pour $\widehat{\Delta^\circ}$ (topos des ens. cosimpliciaux) ? \uuline{\uline{Non}}

Soit $A^{**}\in\operatorname{Ob}\mathbb{DR}_{pd}^{**}(X)$ une Algèbre bigraduée à degrés positifs, alternée pour le degré ext., diff. à p. div.\ldots Supposons que $A^{00}=k$ satisfasse la condition de De Rham, soit $u:k\mathbin{-}\Phi$ d'où $\mathcal{DR}_{pd}^{**}(k,u)$. Considérons alors
\[
\overline A^{**}\unsure{=|.}\mathcal{DR}_{pd}^{**}(A^{**})
\overset{\mathrm{dfn}}=
\mathcal{DR}_{pd}^{**}(k,\unsure{u})\otimes_{k\{T\}}A^{**},
\]
l'augmentation
\[
k\{T\}\mathbin{-}\mathcal{DR}_{pd}^{\unsure{**|*}}(\unsure{k,u})
\]
\inserted{(qui est un quasi-iso \uline{universel})}
définit un hom
\diagram{\begin{tikzcd}
A^{**}\arrow[r]&\overline A^{**}
\end{tikzcd}}
qui est un quasi-iso.
\deleted{\unsure{D'après}}
\inserted{En effet}

\margin{90}
